\scanpage{58}

Let us also put
\[
(90)\qquad
\left\{\begin{aligned}
\nu'&=\nu_3-\nu_2-3\nu_2\nu_3+3\nu_0\\
\nu''&=\nu_3-\nu_2\unsure{\deleted{+}-|\deleted{-}+}2\nu_2\cdot\nu_3+2\nu_0.
\end{aligned}\right.
\]

We then have the relations
\[
(91)\qquad
\left\{\begin{aligned}
\mu&=1-2\nu_2+4\nu'&&({}=\varepsilon_2+4\nu')\\
\mu&=1-2\nu_3+6\nu''&&({}=\varepsilon_3+6\nu'')\\[6pt]
&\cdot\\[-3pt]
\underbrace{-\nu_2+2\nu'}_{\nu}&=
\underbrace{-\nu_3+3\nu''}_{\nu}
&&\text{i.e. }2\nu'-\nu_2+\nu_3=3\nu''.
\end{aligned}\right.
\]
\margin{where $\varepsilon_p=(-1)^{\nu_p}$
\par whence \deleted{$\unsure{\varepsilon_p=1-2\nu_p}$}
\par \deleted{whence} $\varepsilon_p-1=-2\nu_p$
\par $(p\in\{2,3\})$}

\uline{N.B.} The condition $\underline\mu=1$ is equivalent to the conditions
\[
(92)\qquad \nu_2=\nu_3=\nu_0=0
\]
and implies
\[
(93)\qquad \nu=\nu'=\nu''=0
\]

\scanpage{59}
\mspage{677}

4. \unsure{The groups|The group}
\deleted{$\widetilde Z_{\rho,\sigma}$ \unsure{and the sequence} \illegible{2}}
$N_{\rho,\sigma}$, \deleted{\unsure{exact}}\inserted{and the}
exact sequence
\diagram{\begin{tikzcd}[column sep=small]
1 \arrow[r] & N_{\rho,\sigma}^{0+} \arrow[r] & N_{\rho,\sigma} \arrow[r] & \mathbb O_{12}^{*} \arrow[r,dash] & 1.
\end{tikzcd}}
Relations with $\mathrm{Gl}(2,\mathbb Z/12\mathbb Z)$.

We define it as embedded in
$\mathbb O_{12}^{*}\times\underline M\times\underline M$:
\[
\begin{aligned}
(94)\quad&\widetilde Z_{\rho,\sigma}
\subset \mathbb O_{12}^{*}\times\underline M\times\underline M\\[-2pt]
&\prime\prime\\[-2pt]
&\bigl\{(\underline\mu,\alpha_0,\beta_0)\mid
(1-\rho_{\underline M})(\alpha_0)-(1-\sigma_{\underline M})(\beta_0)
=3\nu''(\underline\mu)\lambda_1\bigr\}
\end{aligned}
\]

Note that we have
\[
(95)\qquad (1-\rho_{\underline M})
(\underbrace{\lambda_1-\lambda_0})=3\lambda_1
\]
so the cocycle relation in (\unsure{72|77})
thus takes the form
\[
(96)\qquad
(1-\rho_{\underline M})(\alpha_0-\nu''(\lambda_1-\lambda_0))
-(1-\sigma_{\underline M})(\beta_0)=0
\]
I shall assume \unsure{again} that we have a natural group structure on
$\widetilde Z_{\rho,\sigma}$, for which the projection
$\widetilde Z_{\rho,\sigma}\to\mathbb O_{12}^{*}$
is a group homomorphism. We now find an exact sequence

(97)
\diagram{\begin{tikzcd}[column sep=large]
1 \arrow[r,dash] & \widetilde Z_{\rho,\sigma}^{+0\deleted{\unsure{*}}}
\arrow[r] & \widetilde Z_{\rho,\sigma} \arrow[r] & \mathbb O_{12}^{*} \arrow[r,dash] & 1
\end{tikzcd}}
where
\[
\begin{aligned}
(98)\quad&\widetilde Z_{\rho,\sigma}^{+0}
\overset{\mathrm{dfn}}=\operatorname{Ker}
(\widetilde Z_{\rho,\sigma}\mathbin{-}\mathbb O_{12}^{*})\\[-2pt]
&\rotatebox{90}{$\simeq$}\\[-2pt]
&\bigl\{(\alpha_0,\beta_0)\mid
(1-\rho_{\underline M})(\alpha_0)-(1-\sigma_{\underline M})(\beta_0)=0\bigr\}
\end{aligned}
\]
This group scheme is smooth in characteristic $\ne3$, but not in characteristic $3$, where its dimension goes

\scanpage{60}

from relative dimension $2$ to relative dimension $3$.
This fact of an exceptional fiber of larger dimension cannot, it seems to me, be
eliminated by a dodge, \deleted{\unsure{of}} by introducing
\deleted{some} some ad-hoc definitions for a ``better'' group
$\widetilde Z_{\rho,\sigma}$: there are ``more'' automorphisms of
$\underline{\unsure{\mathcal G}}_k$ (for $k$ an
algebraically closed field), having the properties envisaged, if $k$ is of characteristic $3$, than in the other cases of characteristic.

However, it seems to me that we can construct a subgroup scheme
\deleted{\unsure{of}} of $\widetilde Z_{\rho,\sigma}$, which I feel like denoting
$N_{\rho,\sigma}$, which is smooth over $S_0=\operatorname{Spec}(\mathbb Z)$ and
coincides with $\widetilde Z_{\rho,\sigma}$ over
\deleted{$\unsure{S_{0,3}}$}
$S_{0,\unsure{3|2}}\unsure{{}=}\operatorname{Spec}\mathbb Z[1/3]=S_0\setminus\{3\}$.
To describe it, I note that in the loop equation\inserted{ (96)},
the \deleted{right-hand}\inserted{left-hand} side is
\unsure{automatically|automatic} in
\deleted{$\underline M$, \unsure{which}}\inserted{the submodule of $\underline M$ corresponding}
to the congruence
\[
c_0+c_1\equiv0\ (3)\qquad(\mathrm{cf}\ 62)
\]
and if we introduce the corresponding sub-\deleted{group}\inserted{$\mathbb Z$-module} of index

\scanpage{61}
\mspage{678}

$3$, $M_0$ of $M$
\[
\begin{aligned}
(99)\qquad&M_0\subset M\\[-2pt]
&\prime\prime\ \mathrm{dfn}\\[-2pt]
&\{c_0\lambda_0+c_1\lambda_1\mid c_0+c_1\equiv0\ (3)\},
\end{aligned}
\]
whence a vector bundle $\underline M_0$ over $S_0$ ---
we in fact have a morphism of schemes

(100)
\diagram{\begin{tikzcd}[column sep=large]
\widetilde Z_{\rho,\sigma} \arrow[r] & \underline M_0\\
(\underline\mu,\alpha_0,\beta_0) \arrow[r,mapsto] &
(1-\rho_{\underline M})
(\underbrace{\alpha_0+\nu''(\lambda_1-\lambda_0)}_{\widetilde\alpha_0})
+(1-\sigma_{\underline M})(\beta_0)
\end{tikzcd}}
which we can make explicit as follows: we have a homomorphism of \uline{$\mathbb Z$-modules}

(101)
\diagram{\begin{tikzcd}[column sep=large]
M\times M\times\mathbb Z \arrow[r] & M_0\\
(\alpha_0,\beta_0,\nu'') \arrow[r,mapsto] &
(1-\rho_M)(\alpha_0+\nu''(\lambda_1-\lambda_0))
+(1-\sigma_M)(\unsure{\beta_0})
\end{tikzcd}}
whence a homomorphism \unsure{for} the corresponding module schemes

(102)
\diagram{\begin{tikzcd}[column sep=large]
\underline M\times\underline M\times\underbrace{\underline{\mathbb Z}}_{\mathbb G_a}
\arrow[r] & \underline{M_0}
\end{tikzcd}}
which we compose with

(103)
\diagram{\begin{tikzcd}[column sep=large]
\widetilde Z_{\rho,\sigma} \arrow[r] & \underline M\times\underline M\times\underline{\mathbb Z}\\
(\underline\mu,\alpha_0,\beta_0) \arrow[r,mapsto] & (\alpha_0,\beta_0,\nu''(\mu))
\end{tikzcd}}
We can then consider the cocycle equation (96) as an equation \uline{in}
$\underline M_0$ --- it thus becomes stricter,

\scanpage{62}

whence a closed subscheme of $\widetilde Z_{\rho,\sigma}$,
\[
\begin{aligned}
(104)\quad&N_{\rho,\sigma}\subset\widetilde Z_{\rho,\sigma}\\[-2pt]
&\prime\prime\\[-2pt]
&\{(\underline\mu,\alpha_0,\beta_0)\in\widetilde Z_{\rho,\sigma}\mid
\text{the \deleted{cocycle}\inserted{loop} equation}\\
&\hspace{5cm}\text{(96) valid in }\underline M_0\}.
\end{aligned}
\]
Let us then consider the morphism
\diagram{\begin{tikzcd}
N_{\rho,\sigma} \arrow[r] & \mathbb O_{12}^{*}
\end{tikzcd}}
\margin{N.B. The group $N_{\rho,\sigma}$ should also be denotable by
$\unsure{M|\widetilde N}_{\rho,\sigma}[0]$, unless I am mistaken it
\par \unsure{Or|In} rather it would be
$(\unsure{M}_{\rho,\sigma}^{\unsure{\sharp|b}}/\unsure{M}_{\rho})[0]$
\deleted{\unsure{according to}} \unsure{with} the notations of
$\S\,49$, \deleted{\unsure{2}}\inserted{\uline{$\mathrm{VI}$}}}
It is clear that this morphism is still an epimorphism, the fiber over a
$\underline\mu=(\nu_2,\nu_3,\nu_0)$ being made up of the
$(\alpha_0,\beta_0)\in\Gamma\,\underline M\times\underline M$
satisfying (96) in $M_0$,
\deleted{We \unsure{take the composite homomorphism}}
\inserted{\unsure{is identified|being identified} with the \unsure{scheme} of}
$(\widetilde\alpha_0,\beta_0)$ satisfying the linear equation
\diagram{\begin{tikzcd}
\deleted{M\times M\times\unsure{\mathbb Z}}
\arrow[r,phantom,"{\deleted{\longrightarrow}}" description]
& \deleted{M_0}
\end{tikzcd}}
in $\widetilde\alpha_0,\beta_0$
\[
(105)\qquad
(1-\rho_{\underline M})(\widetilde\alpha_0)
-(1-\sigma_{\underline M})(\beta_0)=0
\qquad\inserted{\text{in }\underline M_0}.
\]
Now the homomorphism of Modules

(106)
\diagram{\begin{tikzcd}[column sep=large]
{} \arrow[d,phantom,"{\Bigg\{}" description] & \underline M\times\underline M \arrow[r] & \underline M_0\\
{} & (\widetilde\alpha_0,\beta_0) \arrow[r,mapsto] &
(1-\rho_{\underline M})(\widetilde\alpha_0)-(1-\sigma_{\underline M})(\beta_0)
\end{tikzcd}}
is a smooth epimorphism --- being associated to a corresponding homomorphism
\diagram{\begin{tikzcd}
M\times M \arrow[r] & M_0
\end{tikzcd}}
which is \unsure{now|manifestly} an epimorphism.
Thus \deleted{$\unsure{\widetilde Z}$} $N_{\rho,\sigma}$ \uline{is smooth} over $S_0$.
\deleted{\unsure{This}} Being a closed subscheme of $\widetilde Z_{\rho,\sigma}$

\scanpage{63}
\mspage{679}

which coincides with it over $S_0\setminus\{3\}$, it
follows that it is the \uline{schematic closure}
\unsure{that|by}
\[
\begin{gathered}
(107)\qquad
N_{\rho,\sigma}=\widetilde Z_{\rho,\sigma}|(\mathrm{Spec}(\mathbb Z)\setminus\{3\}),\\
\text{(schematic closure)}
\end{gathered}
\]
which implies that it is also a subgroup scheme.

Perhaps this is the moment to look a little more closely at the
sub-$\mathbb Z$-module $M_0$ of $M$,
\[
\begin{aligned}
(108)\qquad M_0
&=\mathbb Z(\lambda_1-\lambda_0)+\mathbb Z(3\lambda_1)\\
&=\mathbb Z\underbrace{(\lambda_1-\lambda_0)}_{\delta}
 +\mathbb Z\unsure{(1-\rho_M)}\underbrace{(\lambda_1-\lambda_0)}_{\delta}\\
&=\mathbb Z(\delta\unsure{,|)}\oplus\mathbb Z(1-\rho_M)(\delta)\\
&=\mathbb Z\delta+\mathbb Z\rho_M\delta.
\end{aligned}
\]
Note that $M_0$ is stable under $\rho$, \inserted{$\sigma$}, in
general by construction, for $\forall\ x\in M$, we
\[
\begin{aligned}
(109)\qquad x&\equiv\rho_M(x)\quad(M_0)
 &&\text{(i.e. }(1-\rho_M)(x)\in M_0\text{)}\\
x&\equiv\sigma_M(x)\quad(M_0)
 &&(1-\sigma_M)(x)\in M_0.
\end{aligned}
\]
\deleted{and}

It is of course also stable under $\tau_M=-\mathrm{id}_M$.

Thus $M_0$ is \unsure{normalized} by $\mathcal G$, hence
\inserted{as a subgroup of $\mathcal G$)} normal in
$\mathcal G$.
\deleted{It is}\inserted{\uline{It therefore corresponds to}} a normal
subgroup of $\mathcal G_0=\mathrm{Gl}(2,\mathbb Z)$,

\scanpage{64}

containing $[\Pi_0,\Pi_0]$, of index $3$ in $\Pi_0$.
I presume that it \deleted{is} \uline{contains} the
congruence subgroup relative to $12$ --- i.e. that it corresponds to a
\inserted{normal} subgroup of $\mathrm{Sl}(2,\mathbb Z/12\mathbb Z)$.
It is of index $3.12=36$ in $\mathrm{Sl}(2,\mathbb Z)$,
\deleted{$72$ in $\mathrm{Gl}(2,\mathbb Z)$},
\deleted{\unsure{while}} that
\[
\mathrm{Sl}(2,\mathbb Z/12\mathbb Z)\simeq
\underbrace{\mathrm{Sl}(2,\mathbb Z/4\mathbb Z)}_{48}\times
\underbrace{\mathrm{Sl}(2,\mathbb Z/\unsure{3}\mathbb Z)}_{24}
\]
--- it would therefore be a subgroup \uline{\unsure{of order}}
$32$ of $\mathrm{Sl}(2,\mathbb Z/12\mathbb Z)$,
\inserted{it is obviously the normal subgroup}
generated by the elements \deleted{\illegible{1}} the image of
\deleted{$[\sigma_\infty,\lambda_\infty]$},
\[
\left\{
\begin{array}{l}
(\sigma(\rho)\rho^{-1})^3=[\lambda_0,\lambda_1],\ \text{or }[\sigma_\infty,\lambda_\infty]\\
\text{cf. (7)}\\
\delta\overset{\mathrm{dfn}}{=}\lambda_1\lambda_0^{-1}
\ \deleted{\unsure{\in\Pi_0}}\\
\rho(\delta)
\end{array}
\right.
\]
or again the \inserted{normal} subgroup generated by the single element
\deleted{by the images in $\mathrm{Sl}(2,\mathbb Z/12\mathbb Z)$ of)}
\[
\deleted{\delta=\lambda_1\lambda_0^{-1},\quad
\rho(\delta)=\lambda_\infty\lambda_1^{-1},\quad
\rho^2(\delta)=\lambda_0\lambda_\infty^{-1}}
\]
\[
(110)\qquad
\left\{
\begin{aligned}
\delta&=\lambda_1\lambda_0^{-1}=\ell_1\ell_0^{-1}
=\begin{pmatrix}1&0\\2&1\end{pmatrix}
 \begin{pmatrix}1&+2\\0&1\end{pmatrix}\\
&=\begin{pmatrix}1&+2\\2&5\end{pmatrix}
\end{aligned}
\right.
\]
The question is therefore whether this normal subgroup of
$\mathrm{Sl}(2,\mathbb Z/12\mathbb Z)$ is well and truly of order $32$,
hence of index $36$ --- its order could be larger, its index
smaller. To begin with, what is the order of $\delta$?

\scanpage{65}
\mspage{680}

modulo $12$? We have
\[
\begin{array}{ll}
\deleted{\unsure{(111)}}&
\left\{\begin{aligned}
\delta^2&=\begin{pmatrix}1&2\\2&5\end{pmatrix}
\begin{pmatrix}1&2\\2&5\end{pmatrix}
=\begin{pmatrix}5&12\\12&29\end{pmatrix}
\equiv\begin{pmatrix}5&0\\0&5\end{pmatrix}\quad(12)\\
\delta^4&\equiv1\quad(12),
\end{aligned}\right.\\[-1.5em]
\text{whence}&
\end{array}
\]
that is already a good point! We have \deleted{moreover}
\inserted{even}
\[
(112)\qquad
\diagram{$\left\{\begin{array}{lcl}
\delta\equiv1\quad(2)&&\delta\not\equiv1\ (4),\ \delta\not\equiv1\ (3)\\
&\Downarrow&\\
\delta^2\equiv1\quad(4)&&\delta^2\equiv-1\quad(3)
\end{array}\right.$}
\]


Let us consider
\[
(113)\qquad
\diagram{\begin{tikzcd}[column sep=small]
\mathrm{Sl}(2,\mathbb Z/12\mathbb Z)\simeq
\underbrace{\mathrm{Sl}(2,\mathbb Z/4\mathbb Z)}_{48}\times
\underbrace{\mathrm{Sl}(2,\mathbb Z/3\mathbb Z)}_{24}
\arrow[r,"u\times v"]&
\underbrace{\mathcal G_0^+/\Pi_0}_{12}\times
\underbrace{\mathcal G_3^+}_{3}
\end{tikzcd}}
\]
where \unsure{:}
\[
\diagram{\begin{tikzcd}
u:\mathrm{Sl}(2,\mathbb Z/4\mathbb Z)\arrow[r]&\mathcal G_0^+/\Pi_0
\end{tikzcd}}
\]
is the canonical surjective homomorphism, whose kernel is isomorphic to
$\deleted{\unsure{\mathbb Z}}\mu\times\mu$ --- the three nonzero
elements of the kernel are \unsure{just} $\lambda_0,\lambda_1,\lambda_\infty$
permuted among themselves by $\rho$, and
$\dot\lambda_1\dot\lambda_0^{-1}=\dot\lambda_1\dot\lambda_0
=\dot\lambda_\infty$ and $\ldots$ i.e. \deleted{\unsure{the}}
the $\dot\lambda_i$ are none other than $\dot\delta$ and its permuted images
under $\rho$ --- and
\[
\diagram{\begin{tikzcd}
v:\mathrm{Sl}(2,\mathbb Z/3\mathbb Z)\arrow[r]&\mathcal G_3
\end{tikzcd}}
\]
is the composite homomorphism
\[
\diagram{\begin{tikzcd}[row sep=tiny]
\mathrm{Sl}(2,\mathbb Z/3\mathbb Z)\arrow[r,"\text{epimorphism}"]&
\unsure{A_4|G_4}\arrow[r,"\text{epimorphism}"]
\arrow[d,phantom,"\cap" description]&\mathcal G_3^+\\
&\mathcal G_{\mathbb P^1(\mathbb Z/3\mathbb Z)}&
\end{tikzcd}}
\]
whose kernel is of order $8$. The kernel of $u\times v$ is of order
$4.8=32$, and it obviously contains $\dot\delta$. This shows that
it is indeed the \inserted{normal} subgroup generated by $\delta$
(whose order a priori was \deleted{well and} a \uline{multiple} of $32$), and